\chapter{Линейна регресия}

\section{Модел на линейната регресия}

\begin{defbox}[~(Модел на линейната регресия)]
\[
  Y = \beta_0 + \beta_1 X + \varepsilon,
\]
където $Y$ е \textbf{отклик} (обяснявана величина), $X$ е
\textbf{предиктор} (обясняваща величина), а~$\varepsilon$ е случайна
грешка. За наблюденията $(X_k,Y_k)_{k=1}^n$: $Y_k=\beta_0+\beta_1
X_k+\varepsilon_k$. Навсякъде по-долу приемаме, че не всички $X_k$ са
равни, т.е.\ $\sum_{k=1}^n(X_k-\bar X)^2>0$ (иначе наклонът $\beta_1$
не е идентифицируем от данните).
\end{defbox}

\begin{rembox}[Метод на най-малките квадрати]
Търсим оценки $b_0,b_1$ на $\beta_0,\beta_1$, минимизиращи сумата на
квадратите на грешките:
\[
  \min_{b_0,b_1} \sum_{k=1}^n \big(y_k-b_0-b_1x_k\big)^2 = \min_{b_0,b_1}\sum_{k=1}^n(y_k-\hat y_k)^2, \qquad \hat y_k := b_0+b_1x_k.
\]
\end{rembox}

\begin{thmbox}[]{Теорема}
Решението на горната задача за минимизация е
\[
  b_1 = \frac{\sum_{k=1}^n y_k(x_k-\bar x)}{\sum_{k=1}^n (x_k-\bar x)^2} \quad(\text{наклон на правата}),
\]
\[
  b_0 = \bar y - \bar x\cdot b_1 \quad(\text{пресечна точка с~ординатата}),
\]
където $\bar x=\tfrac1n\sum_k x_k$, $\bar y=\tfrac1n\sum_k y_k$.
\end{thmbox}

\begin{proofbox}
Нека $f(b_0,b_1)=\sum_{k=1}^n(y_k-b_0-b_1x_k)^2$. Приравняваме
частните производни на нула:
\[
  0 = \frac{\partial f}{\partial b_0} = -2\sum_{k=1}^n(y_k-b_0-b_1x_k),
  \qquad
  0 = \frac{\partial f}{\partial b_1} = -2\sum_{k=1}^n(y_k-b_0-b_1x_k)x_k.
\]
От първото уравнение: $n\bar y - nb_0 - b_1 n\bar x = 0 \Rightarrow
b_0=\bar y-b_1\bar x$. Заместваме във второто:
\[
  \sum_{k=1}^n x_ky_k - nb_0\bar x - b_1\sum_{k=1}^n x_k^2 = 0
  \quad\Longrightarrow\quad
  b_1 = \frac{\sum_k x_ky_k-n\bar x\bar y}{\sum_k x_k^2-n\bar x^2} = \frac{\sum_k y_k(x_k-\bar x)}{\sum_k (x_k-\bar x)^2}
\]
(последното преобразувание използва тъждеството $\sum_k(x_k-\bar
x)^2=\sum_k x_k^2-n\bar x^2=\sum_k x_k(x_k-\bar x)$ и~аналогичното
за числителя). \qedbg
\end{proofbox}

\section{Статистически модел и~оценки на параметрите}

\begin{defbox}[~(Статистически модел)]
$Y_k = \beta_0+\beta_1 X_k+\varepsilon_k$, $k=1,\dots,n$, където
$(\varepsilon_j)_{j=1}^n$ са независими в~съвкупност (случайността е
единствено в~грешките $\varepsilon_k$; стойностите $X_k$ са
фиксирани наблюдения), и:
\begin{enumerate}[label=\arabic*)]
  \item $\mathbb E\varepsilon_j=0$;
  \item $D\varepsilon_j=\sigma^2$ за всяко $j$ (\textbf{хомоскедастичност} - еднаква дисперсия на грешката);
  \item $\varepsilon_j\sim N(0,\sigma^2)$.
\end{enumerate}
\end{defbox}

\begin{rembox}[Оценки по метода на най-малките квадрати]
Оценяваме $\beta_0,\beta_1$ чрез метода на най-малките квадрати
(формулите от предната теорема, приложени към случайните $Y_k$):
\[
  \hat\beta_1 = \frac{\sum_{k=1}^n Y_k(X_k-\bar X)}{\sum_{k=1}^n (X_k-\bar X)^2},
  \qquad
  \hat\beta_0 = \bar Y - \bar X\cdot\hat\beta_1.
\]
От модела: $Y_k\sim N(\beta_0+\beta_1X_k,\sigma^2)$, следователно
$\mathbb EY_k=\beta_0+\beta_1X_k$, $DY_k=\sigma^2$ за всяко $k$.
\end{rembox}

\begin{thmbox}[~(неизместеност на \texorpdfstring{$\hat\beta_1$}{бета1-капа})]{Теорема}
$\mathbb E\hat\beta_1 = \beta_1$.
\end{thmbox}

\begin{proofbox}
Нека $A:=\sum_{k=1}^n(X_k-\bar X)^2$. Тогава
\[
  \mathbb E\hat\beta_1 = \frac{\sum_{k=1}^n (\beta_0+\beta_1X_k)(X_k-\bar X)}{A} = \frac{\beta_0\sum_{k=1}^n(X_k-\bar X) + \beta_1\sum_{k=1}^n X_k(X_k-\bar X)}{A}.
\]
Понеже $\sum_k(X_k-\bar X)=\sum_kX_k-n\bar X=0$, а~по тъждеството от
предната теорема $\sum_k X_k(X_k-\bar X)=\sum_k(X_k-\bar X)^2=A$:
\[
  \mathbb E\hat\beta_1 = \frac{0+\beta_1\cdot A}{A} = \beta_1. \qedbg
\]
\end{proofbox}

\begin{thmbox}[~(неизместеност на \texorpdfstring{$\hat\beta_0$}{бета0-капа})]{Теорема}
$\mathbb E\hat\beta_0 = \beta_0$.
\end{thmbox}

\begin{proofbox}
\[
  \mathbb E\hat\beta_0 = \mathbb E\bar Y - \bar X\,\mathbb E\hat\beta_1 = \frac1n\sum_{k=1}^n\mathbb EY_k - \beta_1\bar X = \frac1n\sum_{k=1}^n(\beta_0+\beta_1X_k) - \beta_1\bar X = \beta_0+\beta_1\bar X-\beta_1\bar X = \beta_0. \qedbg
\]
\end{proofbox}

\begin{thmbox}[~(дисперсии на оценките)]{Теорема}
\[
  D\hat\beta_1 = \frac{\sigma^2}{\sum_{k=1}^n(X_k-\bar X)^2},
  \qquad
  D\hat\beta_0 = \sigma^2\left(\frac1n + \frac{\bar X^2}{\sum_{k=1}^n(X_k-\bar X)^2}\right).
\]
Освен това, $\hat\beta_1$ и~$\hat\beta_0$ са линейни комбинации на
независими нормални случайни величини ($Y_k$), следователно самите
те са нормално разпределени:
\[
  \hat\beta_1 \sim N\Big(\beta_1,\ \frac{\sigma^2}{\sum_k(X_k-\bar X)^2}\Big),
  \qquad
  \hat\beta_0 \sim N\Big(\beta_0,\ \sigma^2\Big(\frac1n+\frac{\bar X^2}{\sum_k(X_k-\bar X)^2}\Big)\Big).
\]
\end{thmbox}

\begin{proofbox}
Нека $A=\sum_{k=1}^n(X_k-\bar X)^2$. Понеже $(Y_k)$ са независими,
$D(cY_k)=c^2DY_k$ и~дисперсията на сума от независими е сума от
дисперсиите (Тема~9):
\[
  D\hat\beta_1 = D\left(\frac{\sum_k Y_k(X_k-\bar X)}{A}\right) = \frac1{A^2}\sum_{k=1}^n (X_k-\bar X)^2\,DY_k = \frac{\sigma^2}{A^2}\sum_{k=1}^n(X_k-\bar X)^2 = \frac{\sigma^2}{A}.
\]
За $\hat\beta_0$, записваме го директно като линейна комбинация на
$(Y_k)$:
\[
  \hat\beta_0 = \frac1n\sum_{k=1}^n Y_k - \bar X\cdot\frac{\sum_k Y_k(X_k-\bar X)}{A} = \sum_{k=1}^n Y_k\left(\frac1n - \frac{\bar X(X_k-\bar X)}{A}\right),
\]
и~отново по независимостта на $(Y_k)$:
\[
  D\hat\beta_0 = \sum_{k=1}^n \left(\frac1n-\frac{\bar X(X_k-\bar X)}A\right)^2 DY_k = \sigma^2\left(\frac1n+\frac{\bar X^2}{A}\right)
\]
(след разгръщане на квадрата и~опростяване, ползвайки отново
$\sum_k(X_k-\bar X)=0$ и~$\sum_k(X_k-\bar X)^2=A$). \qedbg
\end{proofbox}

\section{Оценка на дисперсията на грешката}

\begin{rembox}[Идея]
От модела, $\dfrac{Y_k-\beta_0-\beta_1X_k}\sigma\sim N(0,1)$, и~$\sum_{k=1}^n\Big(\dfrac{Y_k-\beta_0-\beta_1X_k}\sigma\Big)^2\sim
\chi^2(n)$. Но $\beta_0,\beta_1$ не са ни известни - заместваме ги с~оценките $\hat\beta_0,\hat\beta_1$, което (аналогично на загубата на
една степен на свобода при замяната на $\mu$ с~$\hat\mu$ в~Тема~15)
намалява броя на степените на свобода с~2 (по броя оценени
параметри):
\[
  \sum_{k=1}^n\left(\frac{Y_k-\hat\beta_0-\hat\beta_1X_k}\sigma\right)^2 \sim \chi^2(n-2).
\]
\end{rembox}

\begin{defbox}[~(Оценка на дисперсията на грешката)]
Използвайки, че за $Z_1,\dots,Z_{n-2}\sim N(0,1)$ независими,
$\tfrac1{n-2}\sum_j Z_j^2 \xrightarrow{\text{п.с.}} \mathbb EZ_1^2=1$
(УЗГЧ, Тема~12), приближаваме $\tfrac1{n-2}\sum_k\big(\tfrac{Y_k-
\hat\beta_0-\hat\beta_1X_k}\sigma\big)^2\approx1$, откъдето
получаваме оценката за дисперсията на грешката:
\[
  \hat\sigma^2 = \frac{\sum_{k=1}^n \big(Y_k-\hat\beta_0-\hat\beta_1X_k\big)^2}{n-2}.
\]
\end{defbox}

\section{Проверка на хипотези за наклона}

\begin{rembox}[Постановка]
Нека искаме да проверим дали наклонът на регресията е равен на
конкретна стойност $b$:
\[
  H_0:\ \beta_1=b \qquad \text{срещу} \qquad H_1:\ \beta_1\neq b
\]
(аналогично се разглеждат едностранните алтернативи $\beta_1<b$,
$\beta_1>b$).
\end{rembox}

\begin{rembox}[Конструкция на теста]
При $H_0$ (т.е.\ $\beta_1=b$), по теоремата за разпределението на
$\hat\beta_1$: $\hat\beta_1\sim N\big(b,\ \sigma^2/A\big)$, $A=
\sum_k(X_k-\bar X)^2$, следователно
\[
  \frac{\hat\beta_1-b}{\sigma/\sqrt A} \sim N(0,1).
\]
Критичната област за ниво на значимост $\alpha$ (двустранен тест) е
\[
  W = \left\{\ |\hat\beta_1-b| \ge q_{1-\alpha/2}\cdot\frac{\sigma}{\sqrt A}\ \right\},
\]
където $q_{1-\alpha/2}$ е $(1-\alpha/2)$-тият квантил на $N(0,1)$ -
отхвърляме $H_0$, ако наблюдаваната разлика между $\hat\beta_1$ и~предполагаемата стойност $b$ е прекалено голяма спрямо стандартното
отклонение на оценката.
\end{rembox}

\begin{rembox}[На практика: неизвестно \texorpdfstring{$\sigma$}{сигма}]
В горната конструкция приехме $\sigma$ за известно. На практика
$\sigma$ обикновено не е известно и~се замества с~оценката
$\hat\sigma$ от предния раздел, а~разпределението на статистиката
$\tfrac{\hat\beta_1-b}{\hat\sigma/\sqrt A}$ вече не е точно $N(0,1)$,
а~$t(n-2)$ - напълно аналогично на прехода от $Z$-тест към $t$-тест
при неизвестна дисперсия за средната стойност (Тема~15).
\end{rembox}
